% ============================================================
% PREAMBULE — Packages et configuration du mémoire AZIZA
% Compilation : xelatex
% ============================================================

% --- Encodage et police ---
\usepackage{fontspec}
\setmainfont{Times New Roman}
\setsansfont{Arial}
\setmonofont{Consolas}

% --- Tableaux ---
\usepackage{array,booktabs,multirow,tabularx,longtable}
\usepackage[table]{xcolor}
\makeatletter
\providecommand\insert@pcolumn{\insert@column}
\makeatother

% --- Langue française ---
\usepackage{polyglossia}
\setdefaultlanguage{french}
\usepackage[autostyle]{csquotes}

% --- Géométrie de la page (normes IFRI) ---
\usepackage[
  a4paper,
  left=3cm,
  right=2.5cm,
  top=2.5cm,
  bottom=2.5cm,
  headheight=15pt,
]{geometry}

% --- Coupures et typographie ---
\usepackage{microtype}

% --- Contrôle des veuves/orphelines ---
\widowpenalty=10000
\clubpenalty=10000
\setlength{\floatsep}{18pt plus 4pt minus 4pt}
\setlength{\textfloatsep}{24pt plus 6pt minus 6pt}

% --- Contrôle des flotteurs ---
\setcounter{totalnumber}{6}
\setcounter{topnumber}{4}
\setcounter{bottomnumber}{4}

% --- Mathématiques ---
\usepackage{amsmath,amssymb,amsfonts}

% --- Graphiques et images ---
\usepackage{graphicx}
\graphicspath{{figures/}{figures/ch4_eda/}{figures/ch4_forme/}
              {figures/ch4_fond/}{figures/ch4_global/}{diagrams/}{plateforme/}}

% --- Listes ---
\usepackage{enumitem}
\setlist{nosep,leftmargin=2em}

% --- Code source ---
\usepackage{listings}
\lstset{
  language=Python,
  basicstyle=\small\ttfamily,
  keywordstyle=\color{blue},
  stringstyle=\color{red},
  commentstyle=\color{green!50!black},
  frame=single,
  numbers=left,
  numberstyle=\tiny,
  breaklines=true,
  showstringspaces=false,
  tabsize=2,
  captionpos=b,
}

% --- Placement forcé des figures ---
\usepackage{float}

% --- Encadrés, théorèmes ---
\usepackage{tcolorbox}
\tcbuselibrary{most}
\usetikzlibrary{positioning}

\newtcolorbox{treebox}[1]{%
  colback=gray!4,
  colframe=black!25,
  coltitle=black!80,
  fonttitle=\bfseries\small,
  title={#1},
  arc=1.5mm,
  left=6mm,
  right=6mm,
  top=2mm,
  bottom=2mm,
  boxrule=0.4pt,
  drop fuzzy shadow,
}

% --- Hyperliens ---
\usepackage[
  colorlinks=true,
  linkcolor=blue,
  citecolor=blue,
  urlcolor=blue,
  hyperfootnotes=false,
]{hyperref}

% --- Bibliographie ---
\usepackage[backend=biber,style=authoryear,sorting=nyt]{biblatex}
\addbibresource{refs.bib}

% --- En-têtes et pieds de page ---
\usepackage{fancyhdr}
\pagestyle{fancy}
\fancyhf{}
\fancyhead[L]{\small\slshape AZIZA — Pipeline NLP hybride}
\fancyhead[R]{\small\thepage}
\renewcommand{\headrulewidth}{0.4pt}

% --- Style des chapitres ---
\makeatletter
\def\@makechapterhead#1{%
  \vspace*{30\p@}%
  {\parindent \z@ \raggedright \normalfont
   \begin{tikzpicture}
     \draw[line width=0.3pt] (2.5cm,0) -- (\textwidth,0);
     \draw[line width=0.3pt] (2.5cm,0) -- (2.5cm,-2em);
     \node[anchor=base west, text=black!50] at (2.8cm,-0.25cm) {%
       \sffamily\footnotesize\bfseries\MakeUppercase{\@chapapp}%
       \hspace{0.5cm}%
       {\rmfamily\huge\bfseries\thechapter}%
     };
   \end{tikzpicture}\par
   \vspace{1.8cm}%
   {\bfseries\LARGE #1}\par\nobreak
   \vspace{0.4cm}%
   \moveright2.5cm\vbox{\hrule height 0.3pt width \dimexpr\textwidth-2.5cm\relax}\par
   \vskip 40\p@
  }}

\def\@makeschapterhead#1{%
  \vspace*{30\p@}%
  {\parindent \z@ \raggedright \normalfont
   \begin{tikzpicture}
     \draw[line width=0.3pt] (2.5cm,0) -- (\textwidth,0);
     \draw[line width=0.3pt] (2.5cm,0) -- (2.5cm,-2em);
   \end{tikzpicture}\par
   \vspace{1.8cm}%
   {\bfseries\LARGE #1}\par\nobreak
   \vspace{0.4cm}%
   \moveright2.5cm\vbox{\hrule height 0.3pt width \dimexpr\textwidth-2.5cm\relax}\par
   \vskip 40\p@
  }}
\makeatother

% --- Espacement interligne ---
\usepackage{setspace}
\onehalfspacing

% --- Annexes ---
\usepackage[title,titletoc]{appendix}

% --- Césure des URL longues ---
\usepackage{xurl}

% --- Notes de bas de page ---
\usepackage[bottom]{footmisc}

% --- Citations en épigraphe ---
\usepackage{epigraph}
\setlength{\epigraphwidth}{0.7\textwidth}

% --- Commandes personnalisées ---
\newcommand{\AZIZA}{\textsc{Aziza}}
\newcommand{\MISTRAL}{\textsc{Mistral}}
\newcommand{\LLM}{\textsc{LLM}}
\newcommand{\IA}{\textsc{IA}}
\newcommand{\TALN}{\textsc{Taln}}
\newcommand{\AES}{\textsc{Aes}}
\newcommand{\STI}{\textsc{Sti}}
